% One-page CV -- MACHINE LEARNING track (AI for maths, neural theorem proving, evaluation,
% and the systems/GPU work that supports it).
% Content is a selection from full-siddharth-bhat.md; style lives in cvstyle.sty.
\documentclass[9pt,a4paper]{extarticle}
\usepackage{cvstyle}

\begin{document}
\cvheader

\section{Summary}
PhD researcher (Cambridge) working at the intersection of machine learning and formal
mathematics: neural proof synthesis, benchmark and evaluation design, and the symbolic
machinery that makes model output checkable.
Built the retrieval-augmented proof-synthesis pipeline for the \fstar{} proof assistant at
Microsoft Research (\textbf{ICSE 2025 best paper}); first-authored work on
\textbf{dataset defects and evaluation failures in Lean theorem-proving benchmarks}
(ICML 2026); showed that a classical symbolic method beats AlphaGeometry on IMO geometry
(NeurIPS 2024 MATH-AI).
Awarded a \href{https://www.renaissancephilanthropy.org/mathbench-towards-evaluating-natural-language-proofs}{Renaissance Philanthropy AI for Maths}
grant (1 of 30 groups funded from 280+ applicants).
Underneath the ML: a decade of performance-critical \cpp{}/CUDA in production compilers ---
\#3 contributor to LLVM's polyhedral loop optimizer, primary author of Lean's LLVM backend.
Fluent in Python/PyTorch, Lean, \cpp{}, CUDA, Haskell, Rust.

\section{Expertise}
\begin{records}
\item \textbf{Neural theorem proving \& RL (PyTorch):} designed and ran the RAG + RL pipeline
that trains models to generate \fstar{} proofs, combining language-model inference with
SMT-backed symbolic search (Microsoft Research; ICSE 2025 best paper); ongoing work on proof
automation for Lean.
\item \textbf{Evaluation \& benchmarks:} first-author study of dataset defects, contamination
and evaluation failures across Lean theorem-proving benchmarks (ICML 2026); established
stronger symbolic baselines for AI geometry systems, showing Wu's method rivals silver
medalists and lifts AlphaGeometry past gold (NeurIPS 2024 MATH-AI); co-developer of the verified
solver that such benchmarks use as a proof checker.
\item \textbf{Representation learning \& NLP:} first-author work modelling word embeddings as
tuples of feature probabilities (RepL4NLP 2020); embedding pipelines (\texttt{word2vec},
\texttt{GloVe}, \texttt{fasttext}) as a researcher and as a TA for \emph{Natural Language:
Applications}.
\item \textbf{Systems for ML (\cpp{}, CUDA, LLVM/MLIR):} \#3 contributor to Polly, LLVM's
polyhedral loop optimizer --- GPU code generation and unified-memory support for
compute-intensive loop nests; stencil and tiling optimization in PolyMage/PLUTO (SC 2017);
MLIR-based compiler backends targeting accelerators.
\end{records}

\section{Selected Publications}
\begin{pubs}
\pub{Faults in Our Formal Benchmarking: Dataset Defects and Evaluation Failures in Lean Theorem Proving}{Pawan Sasanka Ammanamanchi, \me, Stella Biderman}{ICML 2026}
\pub{Towards Neural Synthesis for SMT-Assisted Proof-Oriented Programming}{Saikat Chakraborty, Gabriel Ebner, \me, Sarah Fakhoury, Sakina Fatima, Shuvendu Lahiri, Nikhil Swamy}{ICSE 2025 (best paper)}
\pub{Verifying Wu's Method can Boost Symbolic AI to Rival Silver Medalists and AlphaGeometry to Outperform Gold Medalists at IMO Geometry}{Shiven Sinha, Ameya Prabhu, Ponnurangam Kumaraguru, \me, Matthias Bethge}{NeurIPS 2024 Workshop MATH-AI}
\pub{Word Embeddings as Tuples of Feature Probabilities}{\me, Alok Debnath, Souvik Banerjee, Manish Shrivastava}{RepL4NLP 2020}
\pub{Certified Decision Procedures for Width-Independent Bitvector Predicates}{\me, L\'eo Stefanesco, Chris Hughes, Tobias Grosser}{OOPSLA 2025}
\pub{Interactive Bit Vector Reasoning using Verified Bitblasting}{Henrik B\"oving, \me, Alex Keizer, Abdalrhman Mohamed, L\'eo Stefanesco, Josh Clune, Clark Barrett, Tobias Grosser}{OOPSLA 2025}
\pub{Optimizing Geometric Multigrid Computation using a DSL Approach}{Vinay Vasista, Kumudha KN, \me, Uday Bondhugula}{Supercomputing (SC) 2017}
\end{pubs}
\vspace{1pt}
Full list (14 papers, incl.\ POPL, OOPSLA, CGO, ITP) on
\href{https://scholar.google.com/citations?user=7Irb-OUAAAAJ&hl=en}{Google Scholar}.

\section{Experience}
\role{Jul--Sep '23}{Microsoft Research, Redmond}{Research Intern}{Built a retrieval-augmented RL pipeline for automated theorem proving: PyTorch models generating \fstar{} proofs, paired with SMT-assisted search. ICSE 2025 best paper.}
\role{Sep--Nov '24}{Amazon Web Services, Automated Reasoning Group, Austin}{Research Intern}{Deciding memory non-interference in a Lean-based ARM symbolic simulator; large-scale symbolic execution.}
\role{May--Jul '19}{Tweag.io, Paris}{Research Intern}{Re-implemented portions of the GHC runtime for Asterius, a Haskell$\to$WebAssembly compiler.}
\role{Mar--Dec '17}{ETH Zurich, Scalable Parallel Computing}{Research Intern}{CUDA code generation in Polly for compute-intensive Fortran loop nests.}
\role{May--Jul '16}{IISc Bangalore}{Research Intern}{PolyMage, a DSL compiler for data-parallel stencil computation; tiling patterns, ISL and PLUTO.}

\section{Software}
\proj{Lean 4}{https://github.com/leanprover/lean4/pulls?q=author\%3Abollu}{co-developed \bvdecide{}, the verified bitvector solver (5th at SMT-COMP 2025) used as a proof checker; primary author of Lean's LLVM backend.}
\proj{Polly (LLVM)}{https://polly.llvm.org/}{\#3 contributor to LLVM's polyhedral loop optimizer: 121 commits, 6000 LoC \cpp{}, GPU codegen and unified memory.}
\proj{lean-mlir}{https://github.com/opencompl/lean-mlir}{formal semantics for MLIR in Lean 4, machine-checking compiler rewrites.}
\proj{Lean 4 Metaprogramming Book}{https://github.com/arthurpaulino/lean4-metaprogramming-book}{authored the chapters on tactics and on metaprogramming for embedded DSLs.}

\section{Education \& Awards}
\edu{PhD, Computer Science}{University of Cambridge; transferred from University of Edinburgh}{2022--2026 (writing up)}
\edu{MS / BTech, Computer Science}{IIIT Hyderabad, India (coursework in NLP and deep learning)}{2015--2021}
\vspace{2pt}
\skillline{Awards}{Renaissance Philanthropy AI for Maths grant (1 of 30 from 280+ applicants)
$\cdot$ Best paper, ICSE 2025 $\cdot$ 5th place, SMT-COMP 2025 (\qfbv{})
$\cdot$ Google Summer of Code student '15, mentor '16 and '18}

\end{document}
